% ===== tr-04 块1：第 4 章 pp.1--13；本片 pp.6--9（原文页 6--9）=====
% [原文 p.6]
考虑 $U$ 上由逐纤维积分定义的 1-形式 $\omega_1,\dots,\omega_n$：
\[
\omega_i(\nu)=\int_{X_b\cap\gamma_i}\iota(\nu)\omega,
\]
其中 $\nu$ 是 $B$ 在 $b$ 处的切向量，我们可以把它提升为 $X_b$ 的法向量场。我们有 $\omega_i=f_*(\omega|_{\gamma_i})$，并且由于 $\omega$ 是闭的，$\omega_i$ 也是闭的。于是在 $U$ 上局部地存在函数 $y_1,\dots,y_n$，使得 $dy_i=\omega_i$。此外，这些函数在 $H_1(X_b,\Z)$ 的基的选取与常数范围内是定义良好的。最后，它们给出定义良好的坐标，这可从以下事实推出：由 McLean 定理，$\nu\mapsto\iota(\nu)\omega$ 给出 $\mathcal{T}_{B,b}$ 与 $H^1(X_b,\R)$ 之间的同构。于是 $y_1,\dots,y_n$ 定义了 $B$ 上一个热带仿射结构的局部坐标。
% [Tutor Note: 积分区域是 X_b ∩ γ_i（不是 γ_i），且 ∂/∂y_j 指 B 上坐标切向量在场中的提升；文字层把花体 \mathcal{T}_{B,b} 抽成普通 T_{B,b}，须按页图。]

\para{仿射结构 2。}我们可以对 $\operatorname{Im}\Omega$ 玩同样的把戏：选取子流形
\[
\Gamma_1,\dots,\Gamma_n\subseteq f^{-1}(U),
\]
它们是 $U$ 上的 $n-1$ 维闭链族，并且使得 $\Gamma_1\cap X_b,\dots,\Gamma_n\cap X_b$ 构成 $H_{n-1}(X_b,\Z)$ 的一组基。我们通过 $\lambda_i=-f_*(\operatorname{Im}\Omega|_{\Gamma_i})$ 来定义 $\lambda_i$，或者等价地
\[
\lambda_i(\nu)=-\int_{X_b\cap\Gamma_i}\iota(\nu)\operatorname{Im}\Omega.
\]
同样，$\lambda_1,\dots,\lambda_n$ 是闭 1-形式，局部有 $\lambda_i=d\check y_i$；而 $\check y_1,\dots,\check y_n$ 又是 $B$ 上一个热带仿射结构的仿射坐标。

\para{McLean 度量。}$H^1(X_b,\R)$ 上的 Hodge 度量由
\[
g(\alpha,\beta)=\int_{X_b}\alpha\wedge\ast\beta
\]
给出，其中 $\alpha,\beta$ 是调和 1-形式；它因而诱导出 $B$ 上的一个度量，后者可以写成
\[
g(\nu_1,\nu_2)=-\int_{X_b}\iota(\nu_1)\omega\wedge\iota(\nu_2)\operatorname{Im}\Omega.
\]
% [Tutor Note: McLean 度量公式中两个因子分别是 ι(ν_1)ω 与 ι(ν_2)ImΩ（一个用 ω、一个用 ImΩ），整体带负号；负号与因子顺序都须按页图，文字层常把负号丢掉。]

Hitchin [42] 的一个关键观察是：这些结构由 Legendre 变换相联系：

\begin{proposition}{4.1.2}
设 $y_1,\dots,y_n$ 是 $B$ 上关于由 $\omega$ 诱导的仿射结构的局部仿射坐标。则局部地存在 $B$ 上的函数 $K$，使得
\[
g(\partial/\partial y_i,\partial/\partial y_j)=\partial^2K/\partial y_i\partial y_j.
\]
此外，$\check y_i=\partial K/\partial y_i$ 构成关于由 $\operatorname{Im}\Omega$ 诱导的仿射结构的一组仿射坐标；而若
\[
\check K(\check y_1,\dots,\check y_n)=\sum \check y_iy_i-K(y_1,\dots,y_n)
\]
是 $K$ 的 Legendre 变换，则
% [原文 p.7]
有
\[
\partial^2\check K/\partial\check y_i\partial\check y_j=g(\partial/\partial\check y_i,\partial/\partial\check y_j).
\]
\end{proposition}

\begin{proof}
取如上所述的开邻域 $U$ 上的族 $\gamma_1,\dots,\gamma_n,\Gamma_1,\dots,\Gamma_n$，使这两组基互为 Poincaré 对偶，即对 $b\in U$ 有 $(\gamma_i\cap X_b)\cdot(\Gamma_j\cap X_b)=\delta_{ij}$。设 $\gamma_1^*,\dots,\gamma_n^*$ 与 $\Gamma_1^*,\dots,\Gamma_n^*$ 分别为 $\Gamma(U,R^1f_*\Z)$ 与 $\Gamma(U,R^{n-1}f_*\Z)$ 的对偶基。由 $\gamma_i$ 的选取，我们得到满足 $dy_i=\omega_i$ 的局部坐标 $y_1,\dots,y_n$，于是特别地
\[
\delta_{ij}=\omega_i(\partial/\partial y_j)=\int_{\gamma_i\cap X_b}\iota(\partial/\partial y_j)\omega,
\]
故 $\iota(\partial/\partial y_j)\omega$ 定义了 $H^1(X_b,\R)$ 中的上同调类 $\gamma_j^*$。类似地，令
\[
g_{ij}=-\int_{\Gamma_i\cap X_b}\iota(\partial/\partial y_j)\operatorname{Im}\Omega;
\]
则 $-\iota(\partial/\partial y_j)\operatorname{Im}\Omega$ 定义了 $H^{n-1}(X_b,\R)$ 中的上同调类 $\sum_i g_{ij}\Gamma_i^*$，并且 $\lambda_i=\sum_j g_{ij}dy_j$。于是
\begin{align*}
g(\partial/\partial y_j,\partial/\partial y_k)&=-\int_{X_b}\iota(\partial/\partial y_j)\omega\wedge\iota(\partial/\partial y_k)\operatorname{Im}\Omega\\
&=g_{jk}.
\end{align*}
% [Tutor Note: 指标归属：ι(∂/∂y_j)ω 对应 γ_j^*，−ι(∂/∂y_j)ImΩ 对应 Σ_i g_{ij}Γ_i^*。这里 j 是「求导方向」，i 是「闭链编号」，二者极易抄反；末式右端原文为 g_{jk}，不是 g_{kj}。]
另一方面，设 $\check y_1,\dots,\check y_n$ 是满足 $d\check y_i=\lambda_i$ 的坐标。则
\[
\partial\check y_i/\partial y_j=g_{ij}=g_{ji}=\partial\check y_j/\partial y_i,
\]
故 $\sum \check y_idy_i$ 是闭 1-形式。于是局部地存在函数 $K$，使得 $\partial K/\partial y_i=\check y_i$ 且 $\partial^2K/\partial y_i\partial y_j=g(\partial/\partial y_i,\partial/\partial y_j)$。一个简单的计算便可确认 $\partial\check K/\partial\check y_i=y_i$。另一方面，
\begin{align*}
g(\partial/\partial\check y_i,\partial/\partial\check y_j)&=g\Bigl(\sum_k\frac{\partial y_k}{\partial\check y_i}\frac{\partial}{\partial y_k},\sum_\ell\frac{\partial y_\ell}{\partial\check y_j}\frac{\partial}{\partial y_\ell}\Bigr)\\
&=\sum_{k,\ell}\frac{\partial y_k}{\partial\check y_i}\frac{\partial y_\ell}{\partial\check y_j}g(\partial/\partial y_k,\partial/\partial y_\ell)\\
&=\sum_{k,\ell}\frac{\partial y_k}{\partial\check y_i}\frac{\partial y_\ell}{\partial\check y_j}\frac{\partial\check y_k}{\partial y_\ell}\\
&=\frac{\partial y_j}{\partial\check y_i}=\frac{\partial^2\check K}{\partial\check y_i\partial\check y_j}.
\end{align*}
\proofend
\end{proof}

于是我们引入带有（多值凸函数的）仿射流形的 \emph{Legendre 变换}这一概念。

% [原文 p.8]
\begin{definition}{4.1.3}
设 $B$ 为一个仿射流形。$B$ 上的一个\emph{多值}函数 $K$ 是开覆盖 $\{(U_i,K_i)\}$ 上的一组函数，使得在 $U_i\cap U_j$ 上 $K_i-K_j$ 是仿射线性的。若 Hessian $(\partial^2K_i/\partial y_j\partial y_k)$ 对所有的 $i$、在任何（等价地，所有）仿射坐标系 $y_1,\dots,y_n$ 中都是正定的，则称 $K$ 是\emph{凸}的。
\end{definition}
% [Tutor Note: 多值函数只要求重叠区上 K_i−K_j 为仿射线性（不是常数），这正是 Hessian 能整体一致而 K 本身多值的原因；符号 K 保留「多值」含义，勿加「单值化」解释。]

给定仿射流形与凸多值函数的一个对 $(B,K)$，$(B,K)$ 的 \emph{Legendre 变换}是一个对 $(\check B,\check K)$，其中 $\check B$ 是 $B$ 的底流形上的一个仿射结构，其坐标局部地由 $\check y_i=\partial K/\partial y_i$ 给出；而 $\check K$ 由
\[
\check K_i(\check y_1,\dots,\check y_n)=\sum \check y_jy_j-K_i(y_1,\dots,y_n)
\]
定义。

\begin{exercise}{4.1.4}
验证 $\check K$ 也是凸的，并且 $(\check B,\check K)$ 的 Legendre 变换是 $(B,K)$。
\end{exercise}
% [Tutor Note: 习题 1.4 实际要验证 Legendre 变换是对合：由 ˇy_i = ∂K/∂y_i 出发，两族坐标的 Jacobi 矩阵互逆，对合性原文未给出细节、留作习题。]

奇怪的是，这种带有 Hessian 度量的仿射流形之间的 Legendre 变换，似乎最早出现在一篇早于镜像对称的统计学著作中，见 [2]。

\section{半平坦镜像对称}
让我们暂时忘掉特殊拉格朗日纤维化。取而代之，我们来看看上一节在 $B$ 上找到的结构如何给出镜像对称的一个玩具版本。

\begin{definition}{4.2.1}
设 $B$ 为一个热带仿射流形。
\begin{enumerate}
\item 记 $\Lambda\subseteq\mathcal{T}_B$ 为由 $\partial/\partial y_1,\dots,\partial/\partial y_n$ 局部生成的格局部系统，其中 $y_1,\dots,y_n$ 是局部仿射坐标。由于过渡映射属于 $\R^n\times GL_n(\Z)$，这是定义良好的。置  % [Tutor Note: 原文此处排印作 ×（1300 dpi 裁图复核），按数学含义应为半直积 ⋊；译文严格照原文排印为 ×。]
\[
X(B):=\mathcal{T}_B/\Lambda.
\]
这是 $B$ 上的一个环面丛。此外，$X(B)$ 带有如下局部地定义的复结构。设 $U\subseteq B$ 是一个带有仿射坐标 $y_1,\dots,y_n$ 的开集，于是 $\mathcal{T}_U$ 有坐标函数 $y_1,\dots,y_n,x_1=dy_1,\dots,x_n=dy_n$。则
\[
q_j=e^{2\pi i(x_j+iy_j)}
\]
给出 $\mathcal{T}_U/\Lambda|_U$ 上的一组全纯坐标，且所诱导的复结构与仿射坐标的选取无关。这称为 $X(B)$ 上的\emph{半平坦}复结构。
% [Tutor Note: 注意 x_j = dy_j 是切丛 T_U 的纤维坐标（即 y_j 的对偶坐标），故 q_j 依赖 (x_j, y_j) 两组实坐标；文字层把 2πi(x_j+iy_j) 抽对，但上标里的 i 极易与指标 i 混淆。]
稍后我们会需要它的一个变体：对 $\epsilon>0$，置
\[
X_\epsilon(B):=\mathcal{T}_B/\epsilon\Lambda.
\]
它带有一个复结构，坐标由
\[
q_j=e^{2\pi i(x_j+iy_j)/\epsilon}
\]
给出。
% [原文 p.9]
（稍后我们会看到，极限 $\epsilon\to0$ 对应于一个「大复结构极限」。）
\item 定义 $\check\Lambda\subseteq\mathcal{T}_B^*$ 为由 $dy_1,\dots,dy_n$ 局部生成的格局部系统，其中 $y_1,\dots,y_n$ 是局部仿射坐标。置
\[
\check X(B):=\mathcal{T}_B^*/\check\Lambda.
\]
当然，$\mathcal{T}_B^*$ 带有典范辛结构，而这一辛结构下降到 $\check X(B)$ 上。\proofend
\end{enumerate}
\end{definition}
% [Tutor Note: 定义 2.1 的两项不可对调：X(B)=T_B/Λ 取切丛、带半平坦复结构；ˇX(B)=T_B^*/ˇΛ 取余切丛、带典范辛结构。Λ 由 ∂/∂y_i 生成，ˇΛ 由 dy_i 生成。定义末尾原文有一个孤立 □，此处照录。]

我们记 $f:X(B)\to B$ 与 $\check f:\check X(B)\to B$ 为这两个环面纤维化；它们显然是互相对偶的。

现在再设 $B$ 上还有一个 Hessian 度量 $g$，其局部势函数为 $K$。则下面的命题表明，事实上 $X(B)$ 与 $\check X(B)$ 都成为 Kähler 流形。

\begin{proposition}{4.2.2}
$K\circ f$ 是 $X(B)$ 上的一个（局部）Kähler 势，定义 Kähler 形式 $\omega=2i\partial\bar\partial(K\circ f)$。该度量是 Ricci 平坦的，当且仅当 $K$ 满足实 Monge--Ampère 方程
\[
\det\frac{\partial^2K}{\partial y_i\partial y_j}=\text{常数}.
\]
\end{proposition}

\begin{proof}
局部地取仿射坐标 $(y_j)$ 与复坐标
\[
z_j=\frac{1}{2\pi i}\log q_j=x_j+iy_j,
\]
我们算出 $\omega=2i\partial\bar\partial(K\circ f)=\frac{i}{2}\sum\frac{\partial^2K}{\partial y_j\partial y_k}dz_j\wedge d\bar z_k$，它显然是正的。此外，若 $\Omega=dz_1\wedge\cdots\wedge dz_n$，则 $\omega^n$ 与 $\Omega\wedge\overline{\Omega}$ 成比例，当且仅当 $\det(\partial^2K/\partial y_j\partial y_k)$ 为常数。\proofend
\end{proof}
% [Tutor Note: 系数核对：ω = 2i∂\bar∂(K∘f) 而 ∂\bar∂K = (1/4)Σ K_{jk}dz_j∧d\bar z_k，故合起来系数为 2i·(1/4)=i/2，与原文证明中的 i/2 一致；不要把 2i 与 i/2 抄混。]

我们把这一 Kähler 流形记作 $X(B,K)$。

对偶地我们有

\begin{proposition}{4.2.3}
在 $\mathcal{T}_B^*$ 的局部典范坐标 $y_i,\check x_i$ 中，$\mathcal{T}_B^*$ 上的复坐标函数 $z_j=\check x_j+i\partial K/\partial y_j$ 在 $\check X(B)$ 上诱导出一个定义良好的复结构，而关于该复结构，典范辛形式 $\omega$ 是某个度量的 Kähler 形式。此外，该度量是 Ricci 平坦的，当且仅当 $K$ 满足实 Monge--Ampère 方程
\[
\det\frac{\partial^2K}{\partial y_j\partial y_k}=\text{常数}.
\]
\end{proposition}

\begin{proof}
容易看出，坐标 $y_j$ 的一个仿射线性变换（从而坐标 $\check x_j$ 的相应变换）导致坐标 $z_j$ 的一个线性变换，
